\chapter{Logs, replay and provenance}\label{ch:provenance}
\chapterintro{Every session writes a log of everything that happened in it, by whom, and to which specification. This
chapter describes the session folder, the four record types, the specification hash, the replay that re-runs a log
without any language model, and the tool contract from Python.}

\section{The session folder}\label{sec:folder}
\index{session folder}
The backend writes each workspace session to \code{sessions/}\meta{yyyymmdd-hhmmss}\code{/} in the repository root. A
stand-alone MCP session or a script writes to the directory you give it. The layout is the same:

\begin{table}[h]
\caption{The session folder.}\label{tab:folder}
\begin{fixtable}{@{}>{\raggedright\arraybackslash}p{0.26\linewidth}X@{}}
Path & Contents \\ \midrule
\code{session.jsonl} & The log: one JSON record per line (Section~\ref{sec:records}). \\
\code{fields/} & The final density of every finished solve, \code{rho\_}\meta{job}\code{.npy}, as float16 (2 bytes per
element). \tool{field_stats}, the surface view and the export read them. \\
\code{exports/} & STL files and design reports. \\
\code{renders/} & The PNG pictures made by \tool{render}, named by call number and view. \\
\code{agent/}\meta{run}\code{/} & One folder per assistant run: the run's Codex event stream (\code{work/events.jsonl}),
its standard error, its final message and, in Propose mode, the seed specification and the private session it worked
in. \\
\end{fixtable}
\end{table}

\section{Records}\label{sec:records}
\index{session log!records}
The log has four record types. The examples below are real records from the paper's agent session on brief b06
(\code{results/X3/b06\_motor\_mount/session/session.jsonl}), reformatted and shortened.

\paragraph{\code{session}} opens the log: time, software version, schema version, a session id and the tool profile.
\begin{jsoncode}
{"type": "session", "t": "2026-10-30T12:19:37.880+00:00", "agentto": "0.0.1",
 "schema": "agentto.spec/v0", "id": "f5f023ceb534", "profile": "full"}
\end{jsoncode}

\paragraph{\code{call}} records one tool call: its sequence number, time, actor, tool, arguments, whether it succeeded,
its result (or error), the hash of the specification after the call, and its duration. Binary payloads such as PNG
images are dropped from the log; their hashes are kept.
\begin{jsoncode}
{"type": "call", "seq": 2, "t": "2026-10-30T12:20:00.035+00:00", "actor": "agent",
 "tool": "set_params",
 "args": {"material_name": "Aluminium 6061", "E": 69000, "nu": 0.33, "h": 1,
          "volume_fraction": 0.2},
 "ok": true, "spec_hash": "be59dc4a6544fc93", "dt_s": 0.0001,
 "result": {"changed": ["E", "nu", "material_name", "volume_fraction", "h"]}}
\end{jsoncode}

\paragraph{\code{result}} records a solve that finished after its call returned: the job, its grid and solver, whether
it was warm-started and from which call, iterations, compliance, volume, compute and wall time, the certificate
verdict and its six checks, and the hash of the specification it solved. This one is the cold sketch; note that
\code{spill\_monitored} is false, so it is not certified.
\begin{jsoncode}
{"type": "result", "seq": 25, "job": "39cd686c8342",
 "result": {"status": "converged", "h": 2.0, "grid": [80, 30, 60],
   "n_elem": 144000, "solver": "mg64", "warm_used": false, "iterations": 80,
   "compliance_phys": 3.9057077848391675, "volume": 0.199999999999788,
   "compute_s": 6.29, "wall_s": 6.67, "certified": false,
   "checks": {"converged": true, "state_solve_accepted": true,
              "no_vram_spill": true, "spill_monitored": false,
              "volume_ok": true, "passive_exact": true},
   "release_rung": false, "spec_hash": "c9bd97c53d933b5b"}}
\end{jsoncode}

\paragraph{\code{answer}} records the answer to a clarification question (who answered, which option or none, and any
comment) and, when previews were computed, their summaries.
\begin{jsoncode}
{"type": "answer", "seq": 3, "t": "2026-10-30T12:20:00.065+00:00",
 "answer": {"choice": null, "label": null, "by": "oracle",
            "comment": "None of these is what the brief means."}}
\end{jsoncode}

The actor is \code{user} for workspace edits, \code{agent} for the assistant and attached MCP clients, \code{system} for
a seed, and \code{script} for Python calls unless you pass another name. A workspace call also carries a
\code{meta} field with its origin and a human-readable label (for example ``accept 3 agent cards''); replay ignores it.

\section{The specification hash}\label{sec:hash}
\index{specification hash}
\code{spec\_hash} is the first 16 hexadecimal digits of the SHA-256 of the specification in canonical JSON (sorted keys,
no spaces). A schema-valid specification is first normalized through the model, so default values count and two
documents that differ only in omitted defaults hash the same. Every call records the hash after it; every solve records
the hash it solved; every report records the hash of its specification. This is how the workspace knows that a result
is stale, and how anyone can check later that a released design belongs to a given specification.

The report also records the SHA-256 of \code{session.jsonl} at the moment of export, which ties the release to the full
history that produced it.

\section{Replay without a language model}\label{sec:replay}
\index{replay}
\begin{pycode}
import sys; sys.path.insert(0, "src")
from agentto.tools import replay
print(replay("sessions/20261031-101500/session.jsonl", root="runs/replayed"))
\end{pycode}

Replay starts a fresh session and re-runs every recorded call in order, with the recorded arguments and actors. It
checks:
\begin{itemize}
  \item every call's success flag and the specification hash after it;
  \item every \tool{render} picture's hash and every exported STL's hash;
  \item every clarification, by reusing the recorded answer instead of asking;
  \item every finished solve, which must reproduce its status, iteration count, certificate verdict, warm start and
  compliance, bit for bit.
\end{itemize}
Calls whose outcome depends on timing (\tool{status}, \tool{wait_job}, \tool{cancel}) are skipped, as are solves that
were cancelled or superseded. The result lists the number of calls, the solves replayed, the calls skipped, every
mismatch, and the final specification hash:

\begin{jsoncode}
{"ok": true, "calls": 36, "solves_replayed": 5, "skipped": 3, "mismatches": [],
 "final_spec_hash": "c9bd97c53d933b5b"}
\end{jsoncode}
(the replay record of the b06 session above).

\begin{notebox}[Why the solves reproduce bit for bit]
The upstream solver assembles its coarsest multigrid matrix with atomic scatter-add operations, whose reduction order
varies from run to run, and at the sharpest projection a difference of $10^{-12}$ can flip single elements of the final
design. Sessions use a deterministic coarsest assembly with a fixed reduction order instead. Each solve also records
the coarse-level policy it used, and replay re-runs it with that policy.
\end{notebox}

\begin{evidencebox}
All 15 complete agent sessions of the paper's study X3, including the one that ended in a refusal, replayed with no
mismatch \claim{C1.4}. Replay checks that the recorded decisions reproduce the recorded results. It does not show that
a fresh language-model run would make the same decisions, and it does not check that the decisions matched the
engineer's intent.
\end{evidencebox}

\section{The tool contract from Python}\label{sec:python}
\index{Session@\code{Session}}
Scripts use the same contract. Each \code{call} is validated, applied and logged exactly as a workspace edit or an agent
call would be, and returns \code{\{"ok", "seq", "spec\_hash", "result"\}} or an \code{"error"} in place of the result.

\needspace{19\baselineskip}
\begin{pycode}
import sys; sys.path.insert(0, "src")
from agentto.tools import Session

if __name__ == "__main__":                     # the GPU worker is a spawned process
    s = Session("runs/demo", profile="full")   # writes runs/demo/session.jsonl
    s.call("new_spec", {"name": "bracket", "size": [100, 50, 50]})
    s.call("set_params", {"volume_fraction": 0.3, "cells_short_edge": 20})
    s.call("add_support", {"id": "wall", "region": {"kind": "face", "face": "x-"}})
    s.call("add_load", {"id": "tip", "force": [0, -2000, 0],
                        "region": {"kind": "patch", "face": "x+",
                                   "center": [25, 25], "size": [10, 20]}})
    print(s.call("validate_spec"))
    print(s.call("ladder"))
    print(s.call("solve", {"rung": "sketch"}))  # blocks until the sketch is done
    s.runner.close()                            # stop the GPU worker
\end{pycode}

\begin{warnbox}[A solve needs a complete specification]
\tool{new_spec} creates a draft without a volume fraction or a resolution, and the schema requires both. Until
\tool{set_params} supplies them, \tool{validate_spec} returns schema errors and \tool{solve} is refused with
``spec is not complete/valid''.
\end{warnbox}

A \code{Session} without a directory (\code{Session(None)}) keeps no log and no density fields, which is enough for
building and validating specifications in a test. A script that solves must protect its entry point, as above,
because the GPU worker is started as a spawned process that re-imports the script.
